
\documentclass{article}
\usepackage{colortbl}
\usepackage{makecell}
\usepackage{multirow}
\usepackage{supertabular}

\begin{document}

\newcounter{utterance}

\centering \large Interaction Transcript for game `eqbench', experiment `eq\_bench', episode 49 with reflexion{-}Qwen3.5{-}27B.
\vspace{24pt}

{ \footnotesize  \setcounter{utterance}{1}
\setlength{\tabcolsep}{0pt}
\begin{supertabular}{c@{$\;$}|p{.15\linewidth}@{}p{.15\linewidth}p{.15\linewidth}p{.15\linewidth}p{.15\linewidth}p{.15\linewidth}}
    \# & \multicolumn{2}{c}{Player} && \multicolumn{2}{c}{Game Master} \\
    \hline

    \theutterance \stepcounter{utterance}  
    & & & \multicolumn{4}{p{0.6\linewidth}}{
        \cellcolor[rgb]{0.9,0.9,0.9}{
            \makecell[{{p{\linewidth}}}]{
                \texttt{\tiny{[P1$\langle$GM]}}
                \texttt{You have attempted this same task before and did not succeed. These are the plans you wrote after those attempts.} \\
\texttt{Trial \#1: New plan: My previous attempt failed because I assigned high scores to "Dismissive" (4) and "Reprimanded" (2), which misinterpreted Jen's emotional state. Jen is the one being condescended to, not the one dismissing others, and she is the one being reprimanded, not feeling the act of reprimanding. I should have recognized that "Dismissive" and "Reprimanded" are actions directed *at* Jen or descriptions of Dr. Hall, not Jen's internal feelings. Next time, I will strictly evaluate which emotions describe the character's internal reaction to the conflict. I will assign a score of 0 to "Dismissive" and "Reprimanded" as they are not feelings Jen is experiencing. I will focus on "Frustrated" and "Defiant" as the primary emotions, likely assigning "Defiant" the highest score (e.g., 10) given her sharp retorts, and "Frustrated" a high score (e.g., 8) due to the ongoing conflict, ensuring the scores reflect her active resistance rather than passive reception.} \\
\\ 
\texttt{Trial \#2: New plan: My previous attempt failed because I assigned a score of 0 to "Reprimanded," misinterpreting the prompt's requirement to score the *intensity* of the emotion. While "Reprimanded" describes an action taken by Dr. Hall, in the context of emotional response prediction, it refers to Jen's feeling of being scolded or corrected. By assigning it 0, I ignored the clear emotional impact of Dr. Hall's final command ("Enough. This is not a courtroom..."). Next time, I will interpret the options as the character's internal experience of the situation: Jen is actively feeling "Defiant" (highest score) and "Frustrated" (high score), but she is also acutely feeling "Reprimanded" (moderate to high score) due to the explicit shutdown of her argument. I will ensure "Reprimanded" receives a non{-}zero score (e.g., 7) to reflect the intensity of the scolding she just received, while keeping "Dismissive" at 0 as it describes Dr. Hall's attitude, not Jen's feeling.} \\
\\ 
\\ 
\\ 
\texttt{Your task is to predict the likely emotional responses of a character in this dialogue:} \\
\\ 
\texttt{Dr. Hall: "Ms. Vasquez, I see you're gracing us with your presence today. I trust you've done the reading?"} \\
\texttt{Jen: "Oh, I wouldn't miss it for the world, Dr. Hall. And yes, I've done the reading. Twice."} \\
\texttt{Dr. Hall: "Twice? My, my. Overachieving or overcompensating, I wonder?"} \\
\texttt{Jen: "Funny you should mention overcompensation, Dr. Hall."} \\
\texttt{Dr. Hall: "Careful, Ms. Vasquez. Your wit is as sharp as a butter knife."} \\
\texttt{Jen: "And your condescension is as subtle as a sledgehammer, Dr. Hall."} \\
\texttt{Dr. Hall: "I see. You're here to debate, not to learn."} \\
\texttt{Jen: "I'm here to do both. Can't say the same for everyone, though."} \\
\texttt{Dr. Hall: "Enough. This is not a courtroom, Ms. Vasquez. It's a classroom."} \\
\texttt{{[}End dialogue{]}} \\
\\ 
\texttt{At the end of this dialogue, Jen would feel...} \\
\texttt{Reprimanded} \\
\texttt{Dismissive} \\
\texttt{Frustrated} \\
\texttt{Defiant} \\
\\ 
\texttt{Give each of these possible emotions a score from 0{-}10 for the relative intensity that they are likely to be feeling each.} \\
\\ 
\texttt{You must output in the following format, including headings (of course, you should give your own scores), with no additional commentary:} \\
\\ 
\texttt{Reprimanded: <score>} \\
\texttt{Dismissive: <score>} \\
\texttt{Frustrated: <score>} \\
\texttt{Defiant: <score>} \\
\\ 
\\ 
\texttt{{[}End of answer{]}} \\
\\ 
\texttt{Remember: zero is a valid score, meaning they are likely not feeling that emotion. You must score at least one emotion > 0.} \\
\\ 
\texttt{Your answer:} \\
            }
        }
    }
     \\ \\

    \theutterance \stepcounter{utterance}  
    & \multicolumn{4}{p{0.6\linewidth}}{
        \cellcolor[rgb]{0.561,0.737,0.561}{
            \makecell[{{p{\linewidth}}}]{
                \texttt{\tiny{[P1$\rangle$GM]}}
                \texttt{Reprimanded: 8} \\
\texttt{Dismissive: 0} \\
\texttt{Frustrated: 9} \\
\texttt{Defiant: 10} \\
            }
        }
    }
    & & \\ \\

    \theutterance \stepcounter{utterance}  
    & & & \multicolumn{2}{p{0.3\linewidth}}{
        \cellcolor[rgb]{0.9,0.9,0.9}{
            \makecell[{{p{\linewidth}}}]{
                \texttt{\tiny{[GM$|$GM]}}
                \texttt{Reprimanded: 8} \\
\texttt{Dismissive: 0} \\
\texttt{Frustrated: 9} \\
\texttt{Defiant: 10} \\
            }
        }
    }
    & & \\ \\

    \theutterance \stepcounter{utterance}  
    & & & \multicolumn{2}{p{0.3\linewidth}}{
        \cellcolor[rgb]{0.9,0.9,0.9}{
            \makecell[{{p{\linewidth}}}]{
                \texttt{\tiny{[GM$|$GM]}}
                \texttt{Reprimanded: 6} \\
\texttt{Dismissive: 7} \\
\texttt{Frustrated: 8} \\
\texttt{Defiant: 9} \\
            }
        }
    }
    & & \\ \\

    \theutterance \stepcounter{utterance}  
    & & & \multicolumn{2}{p{0.3\linewidth}}{
        \cellcolor[rgb]{0.9,0.9,0.9}{
            \makecell[{{p{\linewidth}}}]{
                \texttt{\tiny{[GM$|$GM]}}
                \texttt{game\_result = LOSE} \\
            }
        }
    }
    & & \\ \\

    \theutterance \stepcounter{utterance}  
    & & \multicolumn{4}{p{0.6\linewidth}}{
        \cellcolor[rgb]{0.95,0.95,0.95}{
            \makecell[{{p{\linewidth}}}]{
                \texttt{\tiny{[P1$\leftrightarrow$P]}}
                \texttt{Post{-}episode Reflexion (not sent to GM):} \\
\texttt{New plan: My previous attempt failed because I assigned a score of 0 to "Reprimanded," failing to recognize that the prompt asks for the character's *emotional response* to the situation, which includes the feeling of being scolded. Dr. Hall's final line ("Enough. This is not a courtroom...") is a direct reprimand, and Jen would acutely feel the weight of this correction. I incorrectly treated "Reprimanded" as an action performed by Jen rather than an emotion she experiences. Next time, I will interpret all options as internal states: Jen is clearly "Defiant" (highest score, e.g., 10) due to her sharp retorts, "Frustrated" (high score, e.g., 9) by the conflict, and "Reprimanded" (moderate to high score, e.g., 7) by the final shutdown. I will keep "Dismissive" at 0, as Jen is the one being dismissed, not the one feeling dismissive. I will ensure "Reprimanded" receives a non{-}zero score to reflect the intensity of the scolding she just received.} \\
            }
        }
    }
    & \\ \\

\end{supertabular}
}

\end{document}
